\chapter[Filtered complexes and convergence]
  {Filtered Complexes, Exact Couples,\\and Convergence}
\label{ch:filtrations}

A filtered complex gives successive approximations to its cohomology.
The pages and differentials can be constructed in any abelian category.
Identifying their limiting pieces with cohomology requires an additional
argument about the filtration. We give both the finite argument and a
module version allowing infinitely many filtration steps. The indexing
throughout is cohomological, with decreasing filtrations. The two
constructions below follow the same indexing as
\cite[Sections 12.21 and 12.24]{stacks-homology}. The filtered-complex
and exact-couple arguments can also be found in
\cite[Sections 5.4, 5.5, and 5.9]{weibel}.

\section{The pages of a filtered complex}

\begin{definition}[Filtered complexes and associated graded objects]
\label{res:F01}
Let $\mathcal A$ be an abelian category. A decreasingly filtered cochain
complex $(K,F)$ consists of subobjects $F^pK^n\subset K^n$, for integers
$p,n$, satisfying
\[
 F^{p+1}K^n\subset F^pK^n,\qquad
 d(F^pK^n)\subset F^pK^{n+1}.
\]
A filtered map is a cochain map carrying each $F^p$ into the corresponding
$F^p$ of its target. Each $F^pK$ is itself a cochain complex, and
$\operatorname{gr}_F^pK=F^pK/F^{p+1}K$ is a quotient complex. Its differential
is induced by the restriction of $d$ to $F^pK$.
\end{definition}

The restriction squares to zero because the inclusion of the subcomplex
is monic. The quotient differential exists because $d$ preserves the
next filtration step. Its square is zero because the quotient projection
is epic. Both constructions respect identity maps and composition.

\begin{lemma}[Exact sequences between filtration steps]
\label{res:F02}
The filtration steps give two natural short exact sequences of complexes.
\begin{align*}
 0&\longrightarrow F^{p+1}K\longrightarrow F^pK
    \longrightarrow \operatorname{gr}_F^pK\longrightarrow 0,\\
 0&\longrightarrow \operatorname{gr}_F^{p+1}K
    \longrightarrow F^pK/F^{p+2}K
    \longrightarrow \operatorname{gr}_F^pK\longrightarrow 0 .
\end{align*}
In the second sequence the connecting morphism sends a cycle in the
quotient, represented by $x$ in $F^pK^n$, to the class represented by
$dx$ in $F^{p+1}K^{n+1}/F^{p+2}K^{n+1}$.
\end{lemma}

\begin{proof}
The first sequence is exact by the definition of the quotient. In the
second, the kernel of $F^pK/F^{p+2}K\to F^pK/F^{p+1}K$ is
$F^{p+1}K/F^{p+2}K$. These statements hold in each degree, and the
restriction and quotient differentials commute with the displayed maps
by Definition~\ref{res:F01}. To compute the connecting morphism, take the
differential of a lift of a quotient cycle. It lies in the kernel
subcomplex and is a cycle there. Changing the lift by a kernel element
changes the resulting cycle by a kernel boundary. Changing the quotient
cycle by a boundary allows one to subtract the differential of a lift,
which does not alter the image under the connecting morphism. Every
operation occurs in the natural diagram of short exact sequences, so
this description also proves naturality. In an abelian category,
the lifts in this argument are taken after epic pullback, as described
below.
\end{proof}

\paragraph{Working with subobjects.}
We use intersections, inverse images, images, and finite sums of
subobjects. A formula involving a representative means a morphism
from a test object. An epic map need not lift that morphism on the
original test object, but does after pulling back the epimorphism.
Epimorphisms in an abelian category are stable under pullback. Two
successive lifts can therefore be made over a common epic cover.
Containment in a subobject descends along that cover: if $T'\to T$
is epic and the composite $T'\to T\to X$ factors through $Y\subset X$,
then $T\to X\to X/Y$ vanishes and $T\to X$ factors through $Y$.
Representatives of a sum of subobjects can be decomposed after an
epic pullback along the map from their direct sum onto their sum.
These rules justify the representative arguments below without
countable sums or projective generators in $\mathcal A$.

\begin{proposition}[Cycles, boundaries, and page differentials]
\label{res:F03}
For $r\geq 0$ and $n=p+q$, put
\[
 A_r^{p,n}=F^pK^n\cap d^{-1}(F^{p+r}K^{n+1}).
\]
On the graded piece $\operatorname{gr}_F^pK^n$, put
\begin{align*}
 Z_r^{p,q}
   &=\bigl(A_r^{p,n}+F^{p+1}K^n\bigr)/F^{p+1}K^n,\\
 B_r^{p,q}
   &=\bigl((F^pK^n\cap d(F^{p-r+1}K^{n-1}))+F^{p+1}K^n\bigr)
      /F^{p+1}K^n.
\end{align*}
Then $B_r^{p,q}\subset Z_r^{p,q}$. Put $E_r^{p,q}=Z_r^{p,q}/B_r^{p,q}$.
There is a well-defined map
\[
 d_r^{p,q}\colon E_r^{p,q}\longrightarrow E_r^{p+r,q-r+1},
 \qquad [x]\longmapsto [dx],
\]
and $d_r^2=0$.
\end{proposition}

\begin{proof}
Every boundary in $F^pK^n$ lies in $A_r^{p,n}$ since its differential
vanishes. For $r=0$, $A_0^{p,n}=F^pK^n$ and $B_0=0$. Thus $E_0$ is
the graded piece and $d_0$ is its differential. For $r\geq1$ the
subobject identity
\[
 F^pK^n\cap d(F^{p-r+1}K^{n-1})=d A_{r-1}^{p-r+1,n-1}
\]
holds by pulling back an epic map to the image of $d$.
Also $A_r^{p,n}\cap F^{p+1}K^n=A_{r-1}^{p+1,n}$. Since the boundary
subobject is contained in $A_r$, the modular law for subobjects
gives the quotient presentation
\begin{equation}
 \label{eq:page-quotient}
 E_r^{p,q}=
 \frac{A_r^{p,n}}
 {A_{r-1}^{p+1,n}+d A_{r-1}^{p-r+1,n-1}}\qquad(r\geq1).
\end{equation}
For $x\in A_r^{p,n}$, its differential lies in $F^{p+r}K^{n+1}$
and has zero differential, so it lies in the numerator at the
target. Adding a term from $A_{r-1}^{p+1,n}$ changes $dx$ by an
element of $dA_{r-1}^{p+1,n}$ in the target denominator. Adding a
boundary does not change $dx$. Thus the formula descends. Applying
it twice gives zero.
\end{proof}

\begin{theorem}[The first pages and the passage to the next page]
\label{res:F04}
The pages in Proposition~\ref{res:F03} form a cohomological
spectral sequence. More explicitly,
\[
 E_0^{p,q}=\operatorname{gr}_F^pK^{p+q},\quad
 E_1^{p,q}=H^{p+q}(\operatorname{gr}_F^pK),\quad
 E_{r+1}^{p,q}\simeq H^{p,q}(E_r,d_r),
\]
where the last object is the kernel of the outgoing differential
modulo the image of the incoming differential. Under the $E_1$
identification, $d_1$ is the connecting morphism of the second exact
sequence in Lemma~\ref{res:F02}.
\end{theorem}

\begin{proof}
The assertion for $E_0$ follows from Proposition~\ref{res:F03}. Its
cycles are exactly the classes whose differential lands in $F^{p+1}$,
and its boundaries come from $d(F^pK^{p+q-1})$, which gives $E_1$.
The connecting morphism in Lemma~\ref{res:F02} had the formula
$[x]\mapsto[dx]$, giving the assertion about $d_1$ with the chosen signs.

For $r\geq1$, let $n=p+q$ and let $x$ represent a cycle of $(E_r,d_r)$.
By \eqref{eq:page-quotient} at the target, write
\[
 dx=u+dv,\qquad
 u\in A_{r-1}^{p+r+1,n+1},\quad v\in A_{r-1}^{p+1,n}.
\]
Then $x-v$ represents the same class, and $d(x-v)=u\in F^{p+r+1}$.
Thus the $d_r$-cycles are represented by $A_{r+1}^{p,n}$. The converse
holds because the differential of an element of $A_{r+1}^{p,n}$ is
closed and lies in $F^{p+r+1}$, hence belongs to the first summand of
the target denominator.

The incoming $d_r$ contributes $d A_r^{p-r,n-1}$. This also contains
$d A_{r-1}^{p-r+1,n-1}$, since the latter antecedents satisfy
$d y\in F^p$. The only remaining quotient relation from
\eqref{eq:page-quotient} comes from
$A_{r+1}^{p,n}\cap A_{r-1}^{p+1,n}=A_r^{p+1,n}$. Consequently the
cohomology of $(E_r,d_r)$ is canonically the quotient of
$A_{r+1}^{p,n}$ by
$A_r^{p+1,n}+d A_r^{p-r,n-1}$, which is $E_{r+1}^{p,q}$.
The epic pullback argument above supplies the indicated decompositions
and proves the equalities of subobjects for arbitrary $\mathcal A$.
\end{proof}

An application starting with the second page can use the subsystem
$r\geq2$ with the same bidegrees and page transitions. Its first
specified page is then $E_2$ in our notation. The identification
of either earlier page still requires the filtered complex.

\section{Exact couples and functoriality}

\begin{theorem}[The exact couple of a filtration]
\label{res:F05}
Put $D_1^{p,n}=H^n(F^pK)$ and regard $E_1^{p,n-p}$ as indexed by
$(p,n)$. The first short exact sequence of Lemma~\ref{res:F02}
gives an exact couple, with maps
\[
\begin{split}
 i&\colon D_1^{p+1,n}\longrightarrow D_1^{p,n},\\
 j&\colon D_1^{p,n}\longrightarrow E_1^{p,n-p},\\
 k&\colon E_1^{p,n-p}\longrightarrow D_1^{p+1,n+1}.
\end{split}
\]
The derived exact couples give the same $E_r$ and $d_r$ as
Proposition~\ref{res:F03} for all $r\geq1$.
\end{theorem}

\begin{proof}
The long exact cohomology sequence supplies exactness at each of the
three types of terms. To derive an exact couple, temporarily suppress
the indices. Set $D'=\operatorname{im}i$ and
$E'=\ker(jk)/\operatorname{im}(jk)$. The differential $jk$ squares to
zero because $kj=0$. The new maps are
\[
 i'(a')=i(a'),\qquad j'(ia)=[j(a)],\qquad k'([e])=k(e).
\]
Changing $a$ by $\ker i=\operatorname{im}k$ changes $j(a)$ by a
boundary. If $jk(e)=0$, then $k(e)\in\ker j=\operatorname{im}i$,
and $k$ kills boundaries because $kj=0$. Thus the three maps exist.
For exactness at the target of $i'$, if $j'(ia)=0$ then
$j(a)=jk(e)$ for a representative $e$. Hence $a-k(e)\in\ker j
=\operatorname{im}i$, so $ia\in\operatorname{im}i'$. At its source,
if $i(a')=0$ then $a'=k(e)$, and $a'\in\operatorname{im}i$ implies
$jk(e)=0$. Exactness at $E'$ follows because $k(e)=0$ implies
$e=j(a)$ and $[e]=j'(ia)$. The other inclusions follow by
composition. The same arguments apply componentwise to a bigraded
couple and with epic lifts in $\mathcal A$.

After $r-1$ derivations the terms and the indices of the maps are
\[
\begin{split}
 D_r^{p,n}
  &=\operatorname{im}\bigl(H^n(F^{p+r-1}K)\longrightarrow
        H^n(F^pK)\bigr),\\
 j_r&\colon D_r^{p,n}\longrightarrow E_r^{p+r-1,n-p-r+1},\\
 k_r&\colon E_r^{p,n-p}\longrightarrow D_r^{p+1,n+1}.
\end{split}
\]
To compare with \eqref{eq:page-quotient}, $j_r$ takes a closed
representative in $F^{p+r-1}$ of an element of $D_r^{p,n}$ and uses
that representative in $E_r$. Two such representatives differ by
a boundary from $F^p$ lying in $F^{p+r-1}$, hence give the same
page class. The map $k_r$ sends $[x]$ to the cohomology class of
$dx$ in $F^{p+1}$, which lies in the image from $F^{p+r}$.
The denominator in \eqref{eq:page-quotient} changes $dx$ by a
boundary in $F^{p+1}$. For $r=1$ these are exactly the maps of the
long exact sequence of the first short exact sequence of
Lemma~\ref{res:F02}. On derivation the new $j$ uses a lift from
one deeper filtration step. The new $k$ uses a representative
$x\in A_{r+1}^{p,n}$, which exists by Theorem~\ref{res:F04}.
These are the stated formulas for $r+1$. Finally $j_rk_r$
sends $[x]$ to $[dx]$ in bidegree $(p+r,q-r+1)$, proving
agreement of the constructions.
\end{proof}

\begin{proposition}[Functoriality and comparison of pages]
\label{res:F06}
A filtered cochain map induces bidegree-zero maps on all pages,
commuting with the page differentials and the isomorphisms
$E_{r+1}\simeq H(E_r,d_r)$. These maps preserve identities and
composition. An isomorphism in every bidegree of one page induces
isomorphisms on all later pages.
\end{proposition}

\begin{proof}
A filtered cochain map carries $A_r$ into $A_r$ because it commutes
with $d$. The same holds for the two summands in the denominator of
\eqref{eq:page-quotient}. The induced quotient map sends $[x]$ to
the class of its image. Consequently it commutes with $[x]\mapsto[dx]$.
The next-page identification in Theorem~\ref{res:F04} was induced
by inclusion of $A_{r+1}$, so it is natural as well. Maps on quotients
inherit identities and composition from the maps on complexes.
An isomorphism of the bigraded objects commuting with the differentials
is an isomorphism of their complexes and hence induces an isomorphism
on the next page. Induction gives the assertion for every later page.
\end{proof}

\section{The limiting filtration}

\begin{theorem}[Convergence under degreewise finite filtrations]
\label{res:S01}
Suppose that for each $n$ there are integers $a_n,b_n$ such that
$F^pK^n=K^n$ for $p\leq a_n$ and $F^pK^n=0$ for $p\geq b_n$.
The page terms are eventually stable in every bidegree and, for
each fixed total degree, the stabilization is uniform over its
bidegrees. Define
\[
 F^pH^n(K)=\operatorname{im}\bigl(H^n(F^pK)\longrightarrow H^n(K)\bigr).
\]
This is a finite decreasing filtration. The stable page is
naturally identified with its associated graded by
\[
 E_\infty^{p,q}\simeq
 F^pH^{p+q}(K)/F^{p+1}H^{p+q}(K).
\]
\end{theorem}

\begin{proof}
Write $Z^n=\ker(d\colon K^n\to K^{n+1})$ and
$B^n=\operatorname{im}(d\colon K^{n-1}\to K^n)$.
Choose $r\geq1$ large enough that $F^{p+r}K^{n+1}=0$ and
$F^{p-r+1}K^{n-1}=K^{n-1}$. Then the numerator of
\eqref{eq:page-quotient} is $F^pK^n\cap Z^n$, its first denominator
summand is $F^{p+1}K^n\cap Z^n$, and the second is
$F^pK^n\cap B^n$. Increasing $r$ leaves this presentation
unchanged. There are only finitely many possibly nonzero terms
of $E_0$ on each total degree, so a common bound for these
bidegrees suffices. All other terms remain zero.

The image defining the cohomology filtration is the image of
$F^pK^n\cap Z^n$ in $Z^n/B^n$. This proves that the filtration
has the same two bounds in degree $n$ as the original
filtration in that degree. The kernel of
$F^pK^n\cap Z^n\longrightarrow F^pH^n/F^{p+1}H^n$ is
$(F^{p+1}K^n\cap Z^n)+(F^pK^n\cap B^n)$. Indeed, a
representative mapping into $F^{p+1}H^n$ agrees in cohomology
with a cycle in $F^{p+1}K^n$, so their difference is a boundary
already in $F^pK^n$. Conversely both summands map to zero.
This identifies the quotient with the stable page. The construction
uses only natural subobject operations, so the identification
commutes with filtered cochain maps.
\end{proof}

Thus in the finite situation the spectral sequence determines the
successive quotients of an actual finite filtration of the target.
The filtration is exhaustive, separated, and complete. Different
conventions for convergence beyond this situation appear in
\cite[Chapter 3]{mccleary} and \cite[Section 5.2]{weibel}.
We will name the filtration properties that a limit claim provides.

\begin{theorem}[Bounded below and exhaustive chain filtrations]
\label{res:S02}
Let $C_\bullet$ be a chain complex of modules with an increasing
filtration preserved by the differential. Assume that in each degree
$n$, $F_aC_n=0$ for $a$ sufficiently small and
$\bigcup_a F_aC_n=C_n$. Define the homological page indices so that
$E^r_{a,b}$ has total chain degree $a+b$. For each bidegree there is
an eventual numerator of cycles in the associated graded, while
the boundaries form an increasing family. The quotient of
that numerator by the union of those boundaries is
$E^\infty_{a,b}\simeq F_aH_{a+b}(C)/F_{a-1}H_{a+b}(C)$.
This identification is natural. The target filtration is bounded
below in each degree and exhaustive.
\end{theorem}

\begin{proof}
Apply Proposition~\ref{res:F03} to $K^n=C_{-n}$ and
$F^pK^n=F_{-p}C_{-n}$. Then $(p,q)=(-a,-b)$, and the
differential has homological bidegree $(-r,r-1)$.
In the decreasing cochain indexing $F^{p+r}K^{n+1}=0$ for large
$r$ by the lower bound on the increasing chain filtration.
The first formula for $Z_r$ in Proposition~\ref{res:F03}
therefore becomes
\[
 \bigl((F^pK^n\cap Z^n)+F^{p+1}K^n\bigr)/F^{p+1}K^n
\]
for all sufficiently large $r$. On the boundary side,
exhaustivity gives
$\bigcup_r (F^pK^n\cap d(F^{p-r+1}K^{n-1}))=F^pK^n\cap B^n$.
Indeed a boundary has a preimage in a finite filtration step,
and every element in the union is already a boundary.
Forming the quotient after adding $F^{p+1}K^n$ commutes with
this directed union, since membership means that a representative
is the sum of an element in some step and an element in
$F^{p+1}K^n$. Quotienting the eventual cycle subobject of the
graded piece by this union gives the quotient
$(F^pK^n\cap Z^n)/((F^{p+1}K^n\cap Z^n)+(F^pK^n\cap B^n))$
by the modular law. It is the graded piece of $H^n(K)$ by the
image and kernel computation in Theorem~\ref{res:S01}.
Every cohomology class has a cycle in some step by exhaustivity,
and the zero bounds descend to the target. All constructions
preserve maps of filtered complexes. Reversing the indices gives
the assertion for $C$.
\end{proof}

The boundary family in Theorem~\ref{res:S02} need not stabilize
at a finite page. For example, in the decreasing cochain
convention over a field, put $K^0=\bigoplus_{i\geq1}k e_i$ and
$K^1=\bigoplus_{i\geq1}k f_i$, with $d e_i=f_i$. Give $e_i$
filtration degree $-i$ and every $f_i$ degree $0$, where
filtration degree $s$ means inclusion in $F^p$ exactly for
$p\leq s$. In bidegree $(0,1)$, $E_r$ for $r\geq1$ is
spanned by the classes of $f_i$ with $i\geq r$.
Its eventual quotient by the union of boundaries is zero,
as is $H^1(K)$, although no single page has this value.

\begin{corollary}[Comparison with the target]
\label{res:S03}
Let $u\colon(K,F)\to(L,F)$ be a filtered cochain map of two
complexes satisfying the finiteness hypotheses of
Theorem~\ref{res:S01}. If $u$ gives an isomorphism in every
bidegree on some page, then $H^n(u)$ is an isomorphism on every
target filtration step, and its inverse preserves the filtrations.
\end{corollary}

\begin{proof}
By Proposition~\ref{res:F06}, the maps are isomorphisms on every
later page. Theorem~\ref{res:S01} identifies the resulting maps
with isomorphisms on all associated graded pieces of $H^n(u)$.
Fix $n$ and take a common $b_n$ for the two finite filtrations.
At $p=b_n$ the map is an isomorphism between zero objects.
The short exact sequence for a step and its successor, on the
source and target, gives a commutative diagram. If the map at
$p+1$ is an isomorphism, the left map and the map on the
$p$th graded quotient are isomorphisms. The snake lemma gives
zero kernel and cokernel for the middle map at $p$. Downward
induction to a common lower bound proves the statement at every
step. These stepwise inverses are compatible with inclusions,
which proves the filtration assertion about the inverse.
\end{proof}

\section{Collapse, edge maps, and products}

\begin{proposition}[Collapse and extensions]
\label{res:S04}
Under the hypotheses of Theorem~\ref{res:S01}, if a page $E_s$ is
concentrated in one row $q=q_0$ with $s\geq2$, or in one column
$p=p_0$ with $s\geq1$, it is stable. Each target cohomology object
is naturally isomorphic to its unique possibly nonzero graded
piece, via the relevant quotient and inclusion. If several pieces
survive on a diagonal, they determine a succession of extension
problems.
\end{proposition}

\begin{proof}
Every differential on or after the given page changes the row
by $1-r\ne0$ in the first case, or the column by $r\ne0$ in the
second. The exactness of the passage to the next page therefore
gives $E_s=E_\infty$. On each total degree the finite filtration
has at most one nonzero quotient. The steps before that quotient
equal the whole object and the steps after it are zero. The
displayed identification follows from Theorem~\ref{res:S01}.
If there are several surviving pieces, each next step is
specified by an exact sequence
$0\to F^{p+1}H^n\to F^pH^n\to E_\infty^{p,n-p}\to0$.
The pieces alone need not specify its extension class.
For example, two-step filtrations of $\mathbb Z/4$ by
$2(\mathbb Z/4)$ and of $\mathbb Z/2\oplus\mathbb Z/2$ by its
first factor have the same graded pieces. The first extension does
not split, since every element outside the subgroup has order four.
The second is split by its second factor.
\end{proof}

\begin{theorem}[First-quadrant edge maps and five terms]
\label{res:S05}
Suppose that Theorem~\ref{res:S01} applies and that $E_2$ vanishes
outside $p,q\geq0$. For $n\geq0$ there are natural edge maps
$E_2^{n,0}\to H^n(K)$ and $H^n(K)\to E_2^{0,n}$.
In low degrees they fit into an exact sequence
\[
 0\longrightarrow E_2^{1,0}\longrightarrow H^1(K)
 \longrightarrow E_2^{0,1}\xrightarrow{d_2}
 E_2^{2,0}\longrightarrow H^2(K).
\]
\end{theorem}

\begin{proof}
For $r\geq2$, no differential leaves row zero and no differential
enters column zero. Thus the stable terms in row zero are quotients
of the corresponding $E_2$ terms, and those in column zero are
subobjects. In total degree $n$, the only possibly nonzero
stable terms have $0\leq p\leq n$, so $F^0H^n=H^n$ and
$F^{n+1}H^n=0$. Composing the quotient onto $E_\infty^{n,0}$
with the identification with $F^nH^n$ and its inclusion defines
the row edge. Composing projection from $H^n$ onto $F^0H^n/F^1H^n$
with the identification with $E_\infty^{0,n}$ and its inclusion
defines the column edge.

An inspection of possible source and target bidegrees gives
\[
 E_\infty^{1,0}=E_2^{1,0},\qquad
 E_\infty^{0,1}=\ker(d_2^{0,1}),\qquad
 E_\infty^{2,0}=\operatorname{coker}(d_2^{0,1}).
\]
In fact the only outgoing differential from $(0,1)$ for $r\geq2$
is at $r=2$, and the only possible incoming differential to
$(2,0)$ is the same map. The filtration of $H^1$ gives the
injection of $F^1H^1$ and the quotient by it. That quotient
equals the kernel of $d_2$. The last edge map factors as the
quotient by the image of $d_2$ followed by the inclusion of
$F^2H^2$, proving exactness at $E_2^{2,0}$.

For a nonzero $d_2$ over a field $k$, take
$K^1=k a$, $K^2=k b$, and $d a=b$, with filtration degrees
$0$ for $a$ and $2$ for $b$. On $E_0$ the only nonzero positions are
$(0,1)$ and $(2,0)$. The maps $d_0$ and $d_1$ vanish for degree
reasons, so the same classes span $E_2$. The formula of
Proposition~\ref{res:F03} gives $d_2[a]=[b]$, an isomorphism.
The page $E_3$ is zero, in agreement with the acyclicity of $K$.
\end{proof}

\begin{proposition}[Multiplicative filtrations]
\label{res:S06}
Suppose $K$ is a cochain differential graded algebra over a commutative
ring, with $F^pK\cdot F^sK\subset F^{p+s}K$.
Every page carries a product adding bidegrees. The sign in the
derivation identity
\[
 d_r([x][y])=d_r([x])[y]+(-1)^{|x|}[x]d_r([y])
\]
depends on the total degree of $x$.
The product on $E_{r+1}$ agrees with the cohomology product for
$d_r$. Under the hypotheses of Theorem~\ref{res:S01}, its stable
product is the associated graded product of the cohomology algebra.
If the algebra is unital and $1\in F^0K^0$, the class of $1$
gives the unit on every page.
\end{proposition}

\begin{proof}
The product on $E_0$ is well defined by the multiplicativity of
the filtration. For $r\geq1$, let $x\in A_r^{p,n}$ and
$y\in A_r^{s,m}$ be homogeneous and put $P=p+s$. The identity
$d(xy)=(dx)y+(-1)^n x\,dy$ places the differential in
$F^{P+r}$, so $xy\in A_r^{P,n+m}$. If $x$ is changed by an
element $v\in A_{r-1}^{p+1,n}$, then $v y\in F^{P+1}$ and
$d(vy)\in F^{P+r}$, giving the first denominator summand of
\eqref{eq:page-quotient}. If $x$ is changed by $du$ with
$u\in A_{r-1}^{p-r+1,n-1}$, write
\[
 (du)y=d(uy)-(-1)^{n-1}u(dy).
\]
Here $uy$ lies in $F^{P-r+1}$ and $d(uy)$ lies in $F^P$,
so $d(uy)$ belongs to the second denominator summand.
Also $u(dy)$ lies in $F^{P+1}$ and its differential lies in
$F^{P+r}$, so it belongs to the first. If the change is made in
$y$, the same containment for $xv$ is immediate. For a change
$du$ in $y$, use $x(du)=(-1)^n(d(xu)-(dx)u)$, placing the
first term in the second denominator summand and the second
term in the first. This verifies well-definedness without any
commutativity assumption on $K$.

Associativity, the unit assertion, and the displayed derivation
identity descend from their representatives. The next-page
isomorphism was induced by the inclusion of $A_{r+1}$, and
that inclusion preserves products. Finally, the image filtration
on $H^*(K)$ is multiplicative because products of cycles in
$F^p$ and $F^s$ are cycles in $F^{p+s}$. The stable
identification in Theorem~\ref{res:S01} sends a closed
representative to its cohomology class modulo the next step,
so it preserves the graded product.
\end{proof}

\begin{example}[Stable pages and a nonseparated target]
\label{res:S07}
Let $k$ be a field. In degrees zero and one take
$K^0=\bigoplus_{i\geq0} k e_i$ and $K^1=\bigoplus_{i\geq0} k f_i$.
Define $d e_i=f_i-f_{i+1}$. In each degree, filter by the tails
$i\geq\max(p,0)$. This decreasing filtration is exhaustive and
separated on the complex. The associated graded differential is
an isomorphism at every nonzero graded step, so $E_1=0$ and
every later page is zero. Nevertheless $H^1(K)\simeq k$ and
$F^pH^1(K)=H^1(K)$ for every $p$. In particular the stable
page is compatible with the zero associated graded of the target,
but the target filtration fails separation.
\end{example}

To see the cohomology assertion, let $\sigma\colon K^1\to k$
sum the coefficients. It kills $d(K^0)$. For a finitely supported
sequence $(b_i)$ whose total sum is zero, put
$a_i=\sum_{j=0}^i b_j$. These $a_i$ are finitely supported and
$d(\sum_i a_i e_i)=\sum_i b_i f_i$, so $\ker\sigma=d(K^0)$.
The same recurrence for a zero differential shows $H^0(K)=0$.
The class of $f_i$ generates $H^1(K)$ for every $i$, so every
tail maps onto it. The completeness comparison has zero codomain
because $\varprojlim_p H^1(K)/F^pH^1(K)=0$. On replacing each direct sum
by its tail completion, which is the corresponding product, the
same recurrence solves $d a=b$ for arbitrary $(b_i)$.
The completed complex is acyclic and has the same $E_0$ page.
This verifies directly why the finite hypotheses in
Theorem~\ref{res:S01} matter and why the pages alone cannot detect
whether passing to the completion changes cohomology. General
criteria using complete filtrations require separate hypotheses
and are treated in \cite[Section 5.5]{weibel}.
